% ---------------------------------------------------------------------------
% GS2026.2 — Global Solution — Engenharia de Software
% Neighborhood Action Dashboard — CitizenTech
% Formatação segundo ABNT NBR 14724 (estrutura) e NBR 6023 (referências)
% ---------------------------------------------------------------------------
\documentclass[12pt,a4paper,oneside]{article}

\usepackage[T1]{fontenc}
\usepackage[utf8]{inputenc}
\usepackage{times}
\usepackage[top=3cm,bottom=2cm,left=3cm,right=2cm]{geometry}
\usepackage{setspace}
\usepackage{indentfirst}
\usepackage{titlesec}
\usepackage{tocloft}
\usepackage{graphicx}
\usepackage{float}
\usepackage{xcolor}
\usepackage{array}
\usepackage{tabularx}
\usepackage{booktabs}
\usepackage{longtable}
\usepackage{enumitem}
\usepackage{caption}
\usepackage{fancyhdr}
\usepackage{amssymb}
\usepackage{hyperref}
\hypersetup{hidelinks}

\renewcommand{\contentsname}{SUMÁRIO}
\renewcommand{\figurename}{Figura}
\renewcommand{\tablename}{Tabela}
\hyphenpenalty=4000
\tolerance=2500
\emergencystretch=3em

\titleformat{\section}{\normalfont\bfseries\large}{\thesection}{1em}{}
\titleformat{\subsection}{\normalfont\bfseries\normalsize}{\thesubsection}{1em}{}
\titleformat{\subsubsection}{\normalfont\normalsize\itshape}{\thesubsubsection}{1em}{}
\titlespacing*{\section}{0pt}{1.5\baselineskip}{1\baselineskip}
\titlespacing*{\subsection}{0pt}{1.2\baselineskip}{0.8\baselineskip}
\titlespacing*{\subsubsection}{0pt}{1\baselineskip}{0.6\baselineskip}

\renewcommand{\cfttoctitlefont}{\hfill\bfseries\large}
\renewcommand{\cftaftertoctitle}{\hfill}
\setlength{\cftbeforesecskip}{4pt}
\renewcommand{\cftsecfont}{\normalfont}
\renewcommand{\cftsecpagefont}{\normalfont}
\renewcommand{\cftsecleader}{\cftdotfill{\cftdotsep}}

\captionsetup{labelsep=endash,font=small,justification=centering,
              singlelinecheck=false,skip=4pt}
\newcommand{\fonte}[1]{\par\vspace{5pt}{\footnotesize\singlespacing Fonte: #1\par}}

\pagestyle{fancy}
\fancyhf{}
\fancyhead[R]{\small\thepage}
\renewcommand{\headrulewidth}{0pt}
\setlength{\headheight}{14pt}

\newcolumntype{L}[1]{>{\raggedright\arraybackslash}p{#1}}
\newcolumntype{C}[1]{>{\centering\arraybackslash}p{#1}}
\renewcommand{\tabularxcolumn}[1]{>{\raggedright\arraybackslash}p{#1}}
\newcommand{\cab}[1]{\textbf{\footnotesize #1}}
\newcommand{\marcado}{$\checkmark$}
\newcommand{\vazio}{$\square$}

\begin{document}
\onehalfspacing

% =========================================================== CAPA
\thispagestyle{empty}
\begin{center}
{\large\bfseries FACULDADE DE INFORMÁTICA E ADMINISTRAÇÃO PAULISTA — FIAP}\\[2pt]
{\normalsize GRADUAÇÃO EM CIÊNCIA DA COMPUTAÇÃO}\\[2pt]
{\normalsize DISCIPLINA: ENGENHARIA DE SOFTWARE}

\vspace{1.2cm}
{\normalsize
LINCOLN SIMÃO PEREIRA — RM567284\\
MIGUEL SILVA BEZERRA — RM566763\\
NICOLAS SAKAUE NISHIMURA — RM567752\par}

\vspace{2.4cm}
{\normalsize\bfseries GS2026.2 — GLOBAL SOLUTION\par}
\vspace{0.5cm}
{\LARGE\bfseries NEIGHBORHOOD ACTION DASHBOARD\par}
\vspace{0.5cm}
{\large CitizenTech para Gestão Colaborativa\\de Demandas Urbanas\par}
\vspace{0.6cm}
{\normalsize Proposta arquitetural: visão de negócio, arquitetura hexagonal,\\
visões C4, cenário crítico, decisão arquitetural e estratégia de testes\par}

\vfill
{\normalsize SÃO PAULO\\2026\par}
\end{center}
\newpage

% =========================================================== FOLHA DE ROSTO
\thispagestyle{empty}
\begin{center}
{\large\bfseries NEIGHBORHOOD ACTION DASHBOARD\par}
\vspace{0.3cm}
{\normalsize CitizenTech para Gestão Colaborativa de Demandas Urbanas\par}
\end{center}

\vspace{1.2cm}
\begin{flushright}
\begin{minipage}{0.56\textwidth}
\singlespacing\small
Documento apresentado como entrega da Global Solution GS2026.2 da disciplina de
Engenharia de Software, do curso de Ciência da Computação da FIAP, contendo a
concepção arquitetural da plataforma Neighborhood Action Dashboard.

\vspace{6pt}
Prof.\ Dr.\ Leandro Ramos da Silva
\end{minipage}
\end{flushright}

\vspace{1.0cm}
\begin{center}
\begin{minipage}{0.90\textwidth}
\singlespacing\small
\textbf{Integrantes da equipe}\\[6pt]
\begin{tabularx}{\textwidth}{X L{3.2cm}}
\toprule
\cab{Nome completo} & \cab{RM / Matrícula} \\
\midrule
Lincoln Simão Pereira & RM567284 \\
Miguel Silva Bezerra & RM566763 \\
Nicolas Sakaue Nishimura & RM567752 \\
\bottomrule
\end{tabularx}

\vspace{14pt}
\textbf{Repositório do projeto:}
\texttt{https://github.com/lincoln743/GS\_SE}\\[4pt]
\textbf{Abordagem arquitetural adotada:} Arquitetura Hexagonal
(portas e adaptadores)\\[2pt]
\textbf{Visões arquiteturais entregues:} Modelo C4 — diagramas de Contexto e de
Containers\\[2pt]
\textbf{Versão deste documento:} 1.0 — 6 de novembro de 2026
\end{minipage}
\end{center}

\vfill
\begin{center}
{\normalsize SÃO PAULO\\2026\par}
\end{center}
\newpage

% =========================================================== SUMÁRIO
\thispagestyle{empty}
\tableofcontents
\newpage
\setcounter{page}{4}

% =========================================================== 1
\section{INTRODUÇÃO}

\subsection{Apresentação}

A transformação digital ampliou a participação do cidadão na identificação e no
acompanhamento dos problemas que afetam o seu bairro. Iniciativas de
\emph{CitizenTech} usam a tecnologia para aproximar moradores, organizações e
gestores públicos, promovendo transparência, colaboração e decisão apoiada em
dados.

Este documento apresenta a concepção arquitetural do \textbf{Neighborhood Action
Dashboard} (NAD), plataforma digital destinada ao registro, à organização e ao
acompanhamento de demandas de melhoria de bairros e comunidades — falhas de
iluminação pública, buracos em vias, conservação inadequada de praças e parques,
acúmulo de resíduos, problemas de acessibilidade e de sinalização urbana e
necessidades de manutenção em equipamentos públicos.

O foco da entrega não é a implementação: é a definição de uma solução
consistente sob a perspectiva da Engenharia de Software, com visão de negócio,
requisitos, proposta arquitetural, visões, cenário crítico, decisão registrada e
estratégia de qualidade.

\subsection{Escolhas da equipe}

Duas escolhas estruturam o documento e são justificadas nas seções indicadas:

\begin{itemize}[leftmargin=*,itemsep=2pt,topsep=4pt]
  \item \textbf{abordagem arquitetural: Arquitetura Hexagonal} (portas e
        adaptadores) — Seção~\ref{sec:etapa2};
  \item \textbf{visões arquiteturais: Modelo C4}, com os diagramas de Contexto e
        de Containers — Seção~\ref{sec:etapa3}.
\end{itemize}

Conforme a orientação da atividade, apenas uma abordagem arquitetural e apenas
um conjunto de visões foram adotados; as alternativas avaliadas estão
registradas com o motivo do descarte.

\subsection{Guia de leitura}

A Tabela~\ref{tab:guia} relaciona cada etapa solicitada à seção correspondente
deste documento.

\begin{table}[H]
\centering
\caption{Correspondência entre as etapas da Global Solution e as seções}
\label{tab:guia}
{\small\singlespacing
\begin{tabularx}{\textwidth}{L{5.6cm} X L{2.0cm}}
\toprule
\cab{Etapa} & \cab{Seção deste documento} & \cab{Valor} \\
\midrule
Etapa 1 — Visão do Negócio & Seção 2, incluindo requisitos funcionais e não
funcionais & 2,0 \\
Etapa 2 — Arquitetura da Solução & Seção 3, com motivação, benefícios,
limitações e alternativas descartadas & 2,0 \\
Etapa 3 — Visões Arquiteturais & Seção 4, modelo C4: contexto e containers &
2,5 \\
Etapa 4 — Cenário Crítico e Diagrama de Sequência & Seção 5, cenário
\emph{Registrar ocorrência} & 1,5 \\
Etapa 5 — Decisões Arquiteturais (ADR) & Seção 6, ADR-001 & 1,5 \\
Etapa 6 — Qualidade e Testes & Seção 7, pirâmide, tipos de teste e Definition of
Done & 0,5 \\
\bottomrule
\end{tabularx}}
\fonte{elaborado pelos autores (2026), a partir do enunciado da GS2026.2.}
\end{table}

% =========================================================== 2
\section{ETAPA 1 — VISÃO DO NEGÓCIO}
\label{sec:etapa1}

\subsection{Objetivo da solução}

O Neighborhood Action Dashboard tem por finalidade \textbf{centralizar e
padronizar o ciclo de vida das demandas de melhoria urbana de um bairro}, desde
o registro feito pelo cidadão até o encerramento pelo gestor público, oferecendo
a cada parte a informação de que ela precisa: ao cidadão, o andamento do que
registrou; ao gestor, a visão consolidada e priorizada do que precisa ser feito;
à comunidade, a transparência sobre o que foi resolvido e em quanto tempo.

\subsection{Problema a ser resolvido}

Hoje as demandas chegam por canais dispersos — telefone, redes sociais,
atendimento presencial, aplicativos distintos e mensagens informais — e sem
padronização de conteúdo. Três consequências decorrem disso.

A primeira é a \textbf{perda de rastreabilidade}: sem protocolo único, nem o
cidadão sabe o que aconteceu com o seu pedido, nem a gestão consegue dizer
quantas vezes o mesmo problema foi reportado. A segunda é a \textbf{priorização
frágil}: sem dado estruturado sobre categoria, localização, reincidência e
impacto, a fila de atendimento tende a ser ordenada por pressão ou por ordem de
chegada, e não por critério. A terceira é a \textbf{duplicidade}: o mesmo buraco
na mesma rua gera vários registros independentes, que consomem atendimento
repetido e distorcem qualquer indicador.

O problema, portanto, não é a falta de canais — é a ausência de um fluxo único,
padronizado e transparente que una registro, priorização, execução e prestação
de contas.

\subsection{Público-alvo}

A Tabela~\ref{tab:publico} identifica quem usa a solução e quem é beneficiado
por ela.

\begin{table}[H]
\centering
\caption{Público-alvo do Neighborhood Action Dashboard}
\label{tab:publico}
{\small\singlespacing
\begin{tabularx}{\textwidth}{L{3.6cm} L{3.0cm} X}
\toprule
\cab{Grupo} & \cab{Relação} & \cab{Interesse na plataforma} \\
\midrule
Moradores do bairro & Usuário direto & Registrar um problema em menos de dois
minutos e saber, sem telefonar para ninguém, em que pé está \\
Gestores públicos e equipes operacionais & Usuário direto & Receber demandas
padronizadas, priorizadas e georreferenciadas, com histórico \\
Administradores da plataforma & Usuário direto & Manter categorias, usuários e
políticas de priorização sem depender de desenvolvimento \\
Associações de moradores e conselhos locais & Beneficiário & Acompanhar
indicadores do bairro e sustentar a cobrança com dado, não com percepção \\
Comunidade em geral & Beneficiário & Resolução mais ágil, espaço público em
melhor estado e transparência sobre o uso do recurso \\
\bottomrule
\end{tabularx}}
\fonte{elaborado pelos autores (2026).}
\end{table}

\subsection{Benefícios esperados}

\begin{table}[H]
\centering
\caption{Benefícios esperados por grupo de interesse}
\label{tab:beneficios}
{\small\singlespacing
\begin{tabularx}{\textwidth}{X X X}
\toprule
\cab{Para os cidadãos} & \cab{Para os gestores públicos} & \cab{Para a
comunidade} \\
\midrule
Registro simples, com foto e localização capturadas do próprio aparelho &
Visão consolidada de todas as demandas em um único lugar, sem rastrear canais
dispersos & Resolução mais ágil dos problemas do bairro \\
\addlinespace
Protocolo único e acompanhamento transparente, com histórico de situações &
Priorização baseada em dados — categoria, reincidência, impacto e tempo de
espera & Maior participação nas melhorias locais, porque o registro gera
resposta visível \\
\addlinespace
Notificação automática a cada mudança, sem necessidade de cobrar &
Indicadores por bairro e por categoria para dimensionar equipes e orçamento &
Prestação de contas verificável, com dados abertos sobre o que foi resolvido \\
\addlinespace
Menos retrabalho: ocorrências iguais no mesmo ponto são agrupadas &
Menos atendimento duplicado sobre o mesmo problema & Confiança construída sobre
evidência, e não sobre promessa \\
\bottomrule
\end{tabularx}}
\fonte{elaborado pelos autores (2026).}
\end{table}

\subsection{Cenário atual (As-Is) e cenário futuro (To-Be)}

\begin{table}[H]
\centering
\caption{Comparação entre o cenário atual e o cenário futuro}
\label{tab:asis}
{\small\singlespacing
\begin{tabularx}{\textwidth}{L{2.8cm} X X}
\toprule
\cab{Aspecto} & \cab{Cenário atual (As-Is)} & \cab{Cenário futuro (To-Be)} \\
\midrule
Canal de entrada & Telefone, redes sociais, presencial e aplicativos distintos,
sem padrão & Canal único, com formulário padronizado na web e no aplicativo
móvel \\
Identificação & Registro sem protocolo; o cidadão não tem como cobrar &
Protocolo gerado no ato, com consulta pública do andamento \\
Dados da demanda & Texto livre, sem categoria nem localização confiável &
Categoria, descrição, foto e coordenada obrigatórias e validadas \\
Priorização & Por ordem de chegada ou por pressão & Política explícita:
categoria, reincidência no mesmo ponto, impacto e tempo de espera \\
Duplicidade & Cada relato gera um atendimento & Ocorrências próximas da mesma
categoria são agrupadas em uma só \\
Acompanhamento & O cidadão precisa ligar para saber & Notificação automática a
cada mudança de situação \\
Indicadores & Levantamento manual, quando há & Painel por bairro, categoria e
tempo médio de resolução, atualizado continuamente \\
Transparência & Limitada ao que a gestão divulga & Dados abertos exportáveis
sobre demandas e prazos \\
\bottomrule
\end{tabularx}}
\fonte{elaborado pelos autores (2026).}
\end{table}

\subsection{Papéis dos usuários}

\begin{table}[H]
\centering
\caption{Papéis e responsabilidades na plataforma}
\label{tab:papeis}
{\small\singlespacing
\begin{tabularx}{\textwidth}{L{3.0cm} X}
\toprule
\cab{Papel} & \cab{Responsabilidades} \\
\midrule
Cidadão & Registra demandas informando categoria, descrição, localização e
fotos; acompanha o andamento pelo protocolo; confirma ou contesta o
encerramento \\
Gestor Público & Analisa as demandas recebidas, ajusta a prioridade sugerida
pelo sistema, atualiza a situação com justificativa e prazo, e encerra as
solicitações \\
Administrador & Gerencia usuários e perfis de acesso, mantém as categorias de
demanda e configura a política de priorização e os prazos por categoria \\
\bottomrule
\end{tabularx}}
\fonte{elaborado pelos autores (2026), conforme os papéis previstos no enunciado
da atividade.}
\end{table}

\subsection{Requisitos funcionais}

A Tabela~\ref{tab:rf} apresenta os requisitos funcionais da solução, acima do
mínimo de cinco exigido. A prioridade segue a escala Alta, Média e Baixa.

\begin{table}[H]
\centering
\caption{Requisitos funcionais do Neighborhood Action Dashboard}
\label{tab:rf}
{\small\singlespacing
\begin{tabularx}{\textwidth}{L{1.4cm} X L{1.8cm}}
\toprule
\cab{ID} & \cab{Descrição do requisito} & \cab{Prioridade} \\
\midrule
RF01 & Registrar ocorrência informando categoria, descrição, localização
georreferenciada e fotos, gerando protocolo único de acompanhamento & Alta \\
RF02 & Consultar o andamento de uma ocorrência pelo protocolo, exibindo o
histórico completo de situações e as justificativas & Alta \\
RF03 & Classificar e priorizar automaticamente as demandas segundo política
configurável de categoria, reincidência no mesmo ponto, impacto e tempo de
espera & Alta \\
RF04 & Atualizar a situação da ocorrência — Aberta, Em análise, Em execução,
Concluída ou Recusada — com justificativa obrigatória e prazo previsto & Alta \\
RF05 & Exibir painel georreferenciado das demandas, com filtros por bairro,
categoria, situação e período & Alta \\
RF06 & Notificar o cidadão por e-mail ou \emph{push} a cada mudança de situação
da ocorrência que registrou & Média \\
RF07 & Identificar e agrupar ocorrências duplicadas, isto é, da mesma categoria
dentro de um raio configurável e de uma janela de tempo & Média \\
RF08 & Administrar usuários, perfis de acesso, categorias de demanda e prazos
por categoria & Média \\
RF09 & Exportar relatórios e dados abertos em CSV e JSON, com indicadores de
volume, tempo médio de resolução e taxa de reincidência & Baixa \\
\bottomrule
\end{tabularx}}
\fonte{elaborado pelos autores (2026).}
\end{table}

\subsection{Requisitos não funcionais}

A Tabela~\ref{tab:rnf} apresenta os requisitos não funcionais, também acima do
mínimo de quatro exigido, cada um associado ao atributo de qualidade
correspondente.

\begin{table}[H]
\centering
\caption{Requisitos não funcionais e atributos de qualidade}
\label{tab:rnf}
{\small\singlespacing
\begin{tabularx}{\textwidth}{L{1.6cm} X L{3.2cm}}
\toprule
\cab{ID} & \cab{Descrição do requisito} & \cab{Atributo de qualidade} \\
\midrule
RNF01 & O registro de uma ocorrência deve ser concluído em menos de 2 s no
percentil 95, com até 500 usuários simultâneos & Desempenho \\
RNF02 & A plataforma deve manter disponibilidade mensal de 99,5 \%, e o
aplicativo móvel deve permitir registro sem conexão, sincronizando quando a rede
retornar & Disponibilidade e resiliência \\
RNF03 & Dados pessoais do autor da ocorrência não podem ser expostos
publicamente; a plataforma deve atender à LGPD quanto a minimização, finalidade
e retenção, inclusive nas imagens enviadas & Segurança e privacidade \\
RNF04 & As interfaces devem atender às diretrizes WCAG 2.1 no nível AA, e o
registro de uma ocorrência deve ser concluído em no máximo três telas &
Usabilidade e acessibilidade \\
RNF05 & O núcleo de domínio não pode depender de framework web, de banco de
dados ou de serviço externo, de modo que as regras sejam testáveis sem
infraestrutura & Manutenibilidade e testabilidade \\
RNF06 & A plataforma deve suportar a inclusão de novos bairros e o crescimento
do volume de ocorrências sem reprojeto, absorvendo picos por meio de
processamento assíncrono & Escalabilidade \\
RNF07 & As integrações devem ocorrer por API REST documentada em OpenAPI, de
modo que o provedor de mapas, o de identidade e o sistema de protocolo possam
ser substituídos & Interoperabilidade \\
\bottomrule
\end{tabularx}}
\fonte{elaborado pelos autores (2026).}
\end{table}

% =========================================================== 3
\section{ETAPA 2 — ARQUITETURA DA SOLUÇÃO}
\label{sec:etapa2}

\subsection{Abordagem adotada}

\begin{table}[H]
\centering
\caption{Abordagens avaliadas e a escolha da equipe}
\label{tab:escolha}
{\small\singlespacing
\begin{tabularx}{\textwidth}{C{1.6cm} L{4.6cm} X}
\toprule
\cab{Escolha} & \cab{Abordagem} & \cab{Ideia central} \\
\midrule
\vazio & Arquitetura em Camadas & Organização em camadas com responsabilidades
bem definidas \\
\textbf{\marcado} & \textbf{Arquitetura Hexagonal} & \textbf{Núcleo de domínio
isolado, com comunicação por portas e adaptadores} \\
\vazio & Clean Architecture & Dependências apontam para as regras de negócio, no
centro da solução \\
\vazio & Onion Architecture & Domínio ao centro, envolvido por camadas
concêntricas de serviços e infraestrutura \\
\bottomrule
\end{tabularx}}
\fonte{elaborado pelos autores (2026), conforme o quadro da GS2026.2.}
\end{table}

A Figura~\ref{fig:hexagonal} representa a abordagem aplicada ao NAD.

\begin{figure}[H]
\centering
\caption{Arquitetura Hexagonal aplicada ao Neighborhood Action Dashboard}
\label{fig:hexagonal}
\includegraphics[width=\textwidth]{fig_hexagonal.png}
\fonte{elaborado pelos autores (2026), conforme Cockburn (2005).}
\end{figure}

\subsection{Motivação da escolha}

A característica que mais define este problema é a \textbf{multiplicidade de
bordas}. A mesma demanda entra por um portal web, por um aplicativo móvel e, no
futuro próximo, poderá entrar por um canal de mensagens ou pelo atendimento
telefônico transcrito. Ela sai para um provedor de mapas, para um serviço de
notificação, para o sistema de protocolo da prefeitura e para um painel de dados
abertos. Nenhum desses elementos é estável: provedores municipais mudam,
contratos de serviço mudam e canais de atendimento se multiplicam.

A regra de negócio, ao contrário, é estável — o que é uma ocorrência, como ela é
priorizada, quais transições de situação são válidas, quando duas ocorrências
são a mesma. A Arquitetura Hexagonal existe precisamente para separar o que é
estável do que é volátil: o núcleo declara \emph{portas}, interfaces expressas
no vocabulário do domínio, e cada tecnologia de borda entra como um
\emph{adaptador} que as implementa. Acrescentar o canal de WhatsApp é escrever
um novo adaptador primário; trocar o provedor de mapas é escrever um novo
adaptador secundário. Em nenhum dos casos a regra de priorização é tocada.

Há ainda uma razão de prova. O requisito RNF05 exige que as regras sejam
testáveis sem infraestrutura. Com o núcleo isolado por portas, os testes de
unidade da política de priorização e da máquina de situações rodam em
milissegundos, sem subir banco, fila ou servidor web — o que sustenta a base da
pirâmide descrita na Seção~\ref{sec:etapa6}.

\subsection{Benefícios}

\begin{itemize}[leftmargin=*,itemsep=3pt,topsep=4pt]
  \item \textbf{Manutenibilidade}: a mudança de uma tecnologia de borda fica
        confinada a um adaptador; o impacto é local e previsível.
  \item \textbf{Testabilidade}: o domínio é exercitado por testes de unidade
        puros, com dublês nas portas de saída, sem orquestrar infraestrutura.
  \item \textbf{Evolução por canais}: novos pontos de entrada — canal de
        mensagens, integração com ouvidoria, importação em lote — não exigem
        reescrita do núcleo.
  \item \textbf{Substituição de fornecedores}: provedor de mapas, serviço de
        notificação e sistema de protocolo da prefeitura são implementações
        intercambiáveis de portas de saída, o que reduz o aprisionamento
        tecnológico.
  \item \textbf{Clareza de responsabilidade}: a fronteira entre regra de negócio
        e detalhe técnico fica explícita no código, e não apenas na
        documentação.
\end{itemize}

\subsection{Limitações}

Registrar apenas os benefícios seria apresentar a escolha de forma parcial. A
abordagem cobra três preços.

O primeiro é o \textbf{custo estrutural inicial}: portas, adaptadores e objetos
de transporte entre camadas produzem mais arquivos e mais indireção do que uma
organização em camadas simples. Para uma equipe de três integrantes e um prazo
curto, isso é trabalho real.

O segundo é a \textbf{curva de aprendizado}: a inversão de dependência
(``o domínio declara a interface que a infraestrutura implementa'') é
contraintuitiva para quem está habituado ao fluxo direto controlador
$\rightarrow$ serviço $\rightarrow$ repositório, e exige disciplina em revisão
de código para não ser corrompida.

O terceiro é o \textbf{risco de excesso}: aplicada sem critério, a abstração
gera portas com um único implementador que nunca será substituído — custo sem
contrapartida. A equipe adota como regra criar porta apenas quando existir
volatilidade real ou necessidade de dublê em teste; os demais casos usam
chamada direta.

\subsection{Alternativas descartadas}

\begin{table}[H]
\centering
\caption{Alternativas avaliadas e motivo do descarte}
\label{tab:descartadas}
{\small\singlespacing
\begin{tabularx}{\textwidth}{L{3.4cm} X}
\toprule
\cab{Abordagem} & \cab{Por que não foi escolhida} \\
\midrule
Arquitetura em Camadas & É a mais simples de implantar e a mais familiar, mas
não impede que a regra de negócio dependa do framework e do banco: a camada de
serviço costuma importar o ORM diretamente. Como o ponto crítico do NAD é
justamente a troca de integrações externas, a abordagem resolveria o problema
mais fácil e deixaria o mais difícil \\
\addlinespace
Clean Architecture & Atenderia ao problema: o efeito prático sobre a
dependência é equivalente ao da Hexagonal. Foi descartada por acrescentar
camadas prescritivas — entidades, casos de uso, adaptadores de interface e
frameworks — cujo ganho, no tamanho deste sistema, não compensa a cerimônia
adicional. A Hexagonal entrega a mesma inversão com menos estrutura \\
\addlinespace
Onion Architecture & Igualmente próxima em princípio, com o domínio ao centro e
camadas concêntricas. Foi descartada pela mesma razão da Clean, somada à menor
familiaridade da equipe com o seu vocabulário, o que aumentaria o custo de
comunicação sem benefício arquitetural correspondente \\
\bottomrule
\end{tabularx}}
\fonte{elaborado pelos autores (2026).}
\end{table}

Registre-se que Clean e Onion não foram descartadas por serem inadequadas — elas
não são. Foram descartadas por custo de estrutura diante de um ganho
equivalente, e essa é a razão honesta.

\subsection{Respostas às perguntas norteadoras}

\begin{table}[H]
\centering
\caption{Perguntas norteadoras da Etapa 2 e as respostas da equipe}
\label{tab:norteadoras}
{\small\singlespacing
\begin{tabularx}{\textwidth}{L{5.0cm} X}
\toprule
\cab{Pergunta} & \cab{Resposta} \\
\midrule
Quais características de qualidade são mais críticas para o NAD? &
Manutenibilidade e testabilidade, pela volatilidade das integrações municipais;
em seguida, disponibilidade e privacidade, porque o sistema recebe dado pessoal
e imagem de via pública \\
\addlinespace
Como a abordagem separa as regras de negócio da interface, do banco e das
integrações? & O núcleo depende apenas das portas que ele mesmo declara
(\texttt{OcorrenciaRepository}, \texttt{EventoPublisher},
\texttt{GeocodificacaoGateway}, \texttt{ProtocoloGateway}); controladores,
repositórios e clientes externos são implementações dessas interfaces,
resolvidas por injeção de dependência \\
\addlinespace
Que complexidade adicional a abordagem traz, considerando o tamanho da equipe e
o prazo? & Mais arquivos, mais indireção e necessidade de disciplina em revisão.
Mitigação: criar porta apenas onde houver volatilidade ou necessidade de dublê,
e manter um documento curto de convenções no repositório \\
\addlinespace
Por que cada uma das demais abordagens foi descartada? & Camadas, por não
proteger o domínio das integrações; Clean e Onion, por custo de estrutura
superior ao ganho neste porte de sistema — conforme a
Tabela~\ref{tab:descartadas} \\
\addlinespace
Como a escolha se conecta às visões arquiteturais apresentadas? & O núcleo
hexagonal vive dentro do container \emph{API de Demandas} do diagrama de
containers; os adaptadores secundários são exatamente as setas que saem desse
container para o banco, a fila e os sistemas externos
(Seção~\ref{sec:etapa3}) \\
\bottomrule
\end{tabularx}}
\fonte{elaborado pelos autores (2026).}
\end{table}

% =========================================================== 4
\section{ETAPA 3 — VISÕES ARQUITETURAIS}
\label{sec:etapa3}

\subsection{Modelo adotado}

A equipe adota o \textbf{Modelo C4}, proposto por Brown, entregando os diagramas
de \textbf{Contexto} (nível 1) e de \textbf{Containers} (nível 2), conforme
previsto na Opção A da atividade. O Modelo 4+1 de Kruchten não foi utilizado.

A escolha do C4 decorre de dois fatores. O primeiro é a aderência ao que precisa
ser comunicado nesta etapa: o contexto responde \emph{quem usa o sistema e com
que sistemas ele conversa}, e o de containers responde \emph{em que unidades
implantáveis a solução se divide e como elas se comunicam} — exatamente as duas
perguntas que a arquitetura hexagonal escolhida precisa responder para ser
verificável. O segundo é a economia de notação: o C4 se apoia em poucos
elementos e torna o diagrama legível por quem não é da equipe, o que importa em
um sistema cujo interlocutor inclui gestores públicos.

\subsection{Diagrama de Contexto (C4 — nível 1)}

\begin{figure}[H]
\centering
\caption{C4 nível 1 — diagrama de contexto do Neighborhood Action Dashboard}
\label{fig:contexto}
\includegraphics[width=\textwidth]{fig_c4_contexto.png}
\fonte{elaborado pelos autores (2026), conforme Brown (2024).}
\end{figure}

\begin{table}[H]
\centering
\caption{Elementos do diagrama de contexto}
\label{tab:contexto}
{\small\singlespacing
\begin{tabularx}{\textwidth}{L{3.6cm} L{2.2cm} X}
\toprule
\cab{Elemento} & \cab{Tipo} & \cab{Finalidade da interação} \\
\midrule
Cidadão & Pessoa & Registra demandas, informa a localização e acompanha o
andamento pelo protocolo \\
Gestor Público & Pessoa & Analisa, prioriza, atualiza a situação e encerra
solicitações \\
Administrador & Pessoa & Gerencia usuários, categorias e configurações da
plataforma \\
Provedor de Identidade (gov.br) & Sistema externo & Autentica cidadãos e
gestores, dispensando a plataforma de gerir senhas \\
Serviço de Mapas e Geocodificação & Sistema externo & Converte endereço em
coordenada e fornece a base cartográfica do painel \\
Sistema de Protocolo da Prefeitura & Sistema externo & Recebe a demanda e a
integra ao fluxo oficial de atendimento do município \\
Serviço de Notificação & Sistema externo & Entrega e-mail e \emph{push} ao
cidadão a cada mudança de situação \\
Armazenamento de Mídias & Sistema externo & Guarda as fotos anexadas à
ocorrência, fora do banco relacional \\
\bottomrule
\end{tabularx}}
\fonte{elaborado pelos autores (2026).}
\end{table}

\subsection{Diagrama de Containers (C4 — nível 2)}

\begin{figure}[H]
\centering
\caption{C4 nível 2 — diagrama de containers do Neighborhood Action Dashboard}
\label{fig:containers}
\includegraphics[width=\textwidth]{fig_c4_containers.png}
\fonte{elaborado pelos autores (2026), conforme Brown (2024).}
\end{figure}

\begin{table}[H]
\centering
\caption{Containers, tecnologias e responsabilidades}
\label{tab:containers}
{\small\singlespacing
\begin{tabularx}{\textwidth}{L{2.9cm} L{3.0cm} X}
\toprule
\cab{Container} & \cab{Tecnologia} & \cab{Responsabilidade e comunicação} \\
\midrule
Portal Web & React (SPA) & Interface do cidadão e do gestor; consome a API de
Demandas por HTTPS/JSON \\
Aplicativo Móvel & React Native & Registro em campo com foto e coordenada do
aparelho; fila local para operação sem rede; mesmo contrato REST do portal \\
API de Demandas & Node.js com NestJS & Hospeda os adaptadores de entrada e o
núcleo de domínio hexagonal: registro, priorização, transição de situação e
consulta. Acessa o banco por SQL/TLS e publica eventos na fila \\
Serviço de Notificações & Worker Node.js & Consome o evento
\texttt{OcorrenciaRegistrada} e os eventos de mudança de situação e aciona o
serviço externo de notificação \\
Agregador de Indicadores & Worker Node.js & Consolida o painel por bairro,
categoria e tempo de resolução a partir dos eventos, sem onerar a API \\
Banco de Dados & PostgreSQL com PostGIS & Persiste ocorrências, usuários e
categorias; responde às consultas por área, que são a base do painel
georreferenciado \\
Fila de Mensagens & RabbitMQ & Desacopla o registro das notificações e dos
indicadores, absorvendo picos de demanda \\
\bottomrule
\end{tabularx}}
\fonte{elaborado pelos autores (2026).}
\end{table}

\subsection{Alinhamento entre as visões e os requisitos}

A Tabela~\ref{tab:alinhamento} liga cada container aos requisitos que ele
atende, para tornar verificável a afirmação de que a arquitetura responde ao que
foi especificado na Etapa 1.

\begin{table}[H]
\centering
\caption{Rastreabilidade entre containers e requisitos}
\label{tab:alinhamento}
{\small\singlespacing
\begin{tabularx}{\textwidth}{L{3.2cm} L{3.4cm} X}
\toprule
\cab{Container} & \cab{Requisitos atendidos} & \cab{Observação} \\
\midrule
Portal Web & RF01, RF02, RF04, RF05, RF08; RNF04 & Concentra a exigência de
acessibilidade WCAG 2.1 AA \\
Aplicativo Móvel & RF01, RF02; RNF02 & A fila local sustenta o registro sem
conexão \\
API de Demandas & RF01 a RF08; RNF01, RNF05, RNF07 & Abriga o núcleo de
domínio; é onde a política de priorização e o agrupamento de duplicadas vivem \\
Serviço de Notificações & RF06 & Assíncrono, para não atrasar a resposta do
registro \\
Agregador de Indicadores & RF05, RF09 & Pré-consolida os indicadores do painel
e da exportação \\
Banco de Dados & RF02, RF05, RF07; RNF03 & PostGIS viabiliza o raio de
agrupamento e o filtro por bairro \\
Fila de Mensagens & RF06, RF07; RNF06 & Absorve picos e isola falhas das
integrações externas \\
\bottomrule
\end{tabularx}}
\fonte{elaborado pelos autores (2026).}
\end{table}

% =========================================================== 5
\section{ETAPA 4 — CENÁRIO CRÍTICO E DIAGRAMA DE SEQUÊNCIA}
\label{sec:etapa4}

\subsection{Cenário escolhido e justificativa}

O cenário crítico modelado é \textbf{Registrar ocorrência}. A escolha se
justifica por três razões: é o caso de uso de maior volume, pois todo o restante
do sistema depende de ele existir; é o que atravessa o maior número de
participantes, da interface ao domínio, ao banco e à mensageria; e é o único em
que a falha é diretamente visível para o cidadão — se o registro se perde, a
plataforma perde o seu propósito e a confiança que pretende construir.

\subsection{Participantes}

\begin{table}[H]
\centering
\caption{Participantes do cenário crítico}
\label{tab:participantes}
{\small\singlespacing
\begin{tabularx}{\textwidth}{L{4.0cm} X}
\toprule
\cab{Participante} & \cab{Papel no cenário} \\
\midrule
Cidadão & Ator que inicia o fluxo preenchendo o formulário \\
Portal Web & Interface que coleta os dados e trata o retorno \\
API de Demandas & Adaptador primário REST; recebe a requisição e a traduz para
a porta de entrada do domínio \\
Serviço de Demandas & Núcleo de domínio; valida os dados e aplica as regras de
registro \\
Repositório & Adaptador secundário de persistência sobre PostgreSQL \\
Publicador & Adaptador secundário de mensageria sobre RabbitMQ \\
\bottomrule
\end{tabularx}}
\fonte{elaborado pelos autores (2026).}
\end{table}

\begin{figure}[H]
\centering
\caption{Diagrama de sequência do cenário Registrar ocorrência}
\label{fig:sequencia}
\includegraphics[width=\textwidth]{fig_sequencia.png}
\fonte{elaborado pelos autores (2026), conforme Object Management Group (2017).}
\end{figure}

\subsection{Fluxo principal}

\begin{table}[H]
\centering
\caption{Mensagens do fluxo principal — Registrar ocorrência}
\label{tab:fluxo}
{\small\singlespacing
\begin{tabularx}{\textwidth}{C{0.9cm} L{3.0cm} L{3.0cm} X}
\toprule
\cab{\#} & \cab{De} & \cab{Para} & \cab{Mensagem} \\
\midrule
1 & Cidadão & Portal Web & Preenche o formulário (categoria, descrição,
localização, foto) e envia \\
2 & Portal Web & API de Demandas & \texttt{POST /ocorrencias} com os dados da
ocorrência \\
3 & API de Demandas & Serviço de Demandas & \texttt{registrarOcorrencia(dados)}
pela porta de entrada \\
4 & Serviço de Demandas & Serviço de Demandas & \texttt{validar(dados)}:
campos obrigatórios, categoria existente e coordenada dentro da área atendida \\
5 & Serviço de Demandas & Repositório & \texttt{salvar(ocorrência, situação =
ABERTA)} \\
6 & Repositório & Serviço de Demandas & Retorna o protocolo gerado \\
7 & Serviço de Demandas & Publicador & \texttt{publicar(OcorrenciaRegistrada)}
para notificação e indicadores \\
8 & Serviço de Demandas & API de Demandas & Retorna o protocolo \\
9 & API de Demandas & Portal Web & \texttt{201 Created} com o protocolo \\
10 & Portal Web & Cidadão & Exibe o protocolo de acompanhamento \\
\bottomrule
\end{tabularx}}
\fonte{elaborado pelos autores (2026).}
\end{table}

\subsection{Fluxos alternativos}

\begin{table}[H]
\centering
\caption{Fluxos alternativos do cenário}
\label{tab:alternativos}
{\small\singlespacing
\begin{tabularx}{\textwidth}{C{1.0cm} L{3.0cm} L{3.0cm} X}
\toprule
\cab{\#} & \cab{De} & \cab{Para} & \cab{Mensagem} \\
\midrule
A1 & Serviço de Demandas & Portal Web (via API) & [dados inválidos] Retorna
\texttt{422} com o erro por campo; a interface destaca os campos e solicita
correção, preservando o que já foi digitado \\
\addlinespace
A2 & Repositório & Portal Web (via API) & [falha ao salvar] Retorna
\texttt{503}; a interface informa a indisponibilidade, mantém o rascunho local e
oferece novo envio, de modo que nenhum dado do cidadão é perdido \\
\addlinespace
A3 & Serviço de Demandas & Portal Web (via API) & [ocorrência duplicada] Quando
já existe demanda aberta da mesma categoria no raio configurado, o sistema
vincula o novo relato à ocorrência existente e devolve o protocolo dela,
informando o cidadão de que o problema já está em acompanhamento \\
\bottomrule
\end{tabularx}}
\fonte{elaborado pelos autores (2026).}
\end{table}

O fluxo A3 merece nota, por ser decisão de projeto e não apenas tratamento de
erro: ele implementa o requisito RF07 no momento do registro, e não em uma
rotina posterior de limpeza. Evitar a duplicata na entrada custa uma consulta
geoespacial; removê-la depois custa atendimento repetido e indicador distorcido.

% =========================================================== 6
\section{ETAPA 5 — DECISÃO ARQUITETURAL (ADR)}
\label{sec:etapa5}

\begin{table}[H]
\centering
\caption{ADR-001 — Adoção da Arquitetura Hexagonal no núcleo de demandas}
\label{tab:adr}
{\small\singlespacing
\begin{tabularx}{\textwidth}{L{3.4cm} X}
\toprule
\cab{Status / Data} & \textbf{Aceita} \quad (\vazio\ Proposta \quad
\textbf{\marcado}\ Aceita \quad \vazio\ Substituída) \quad Data: 06/11/2026 \\
\midrule
Contexto & O Neighborhood Action Dashboard recebe demandas por múltiplos canais
— portal web e aplicativo móvel hoje, canal de mensagens e importação da
ouvidoria em estudo — e depende de quatro integrações externas sobre as quais a
equipe não tem governança: provedor de identidade, serviço de mapas, sistema de
protocolo da prefeitura e serviço de notificação. Contratos municipais são
revistos periodicamente, e a substituição de qualquer um desses fornecedores é
provável dentro do horizonte do projeto. Ao mesmo tempo, as regras de negócio —
política de priorização, máquina de situações e critério de duplicidade — são
estáveis e precisam ser testadas sem infraestrutura, conforme o RNF05. A equipe
tem três integrantes e prazo curto \\
\addlinespace
Decisão & Decidimos adotar a Arquitetura Hexagonal (portas e adaptadores) no
container \emph{API de Demandas}. O núcleo de domínio declara as portas de
entrada (\texttt{RegistrarOcorrencia}, \texttt{AtualizarSituacao},
\texttt{PriorizarDemanda}, \texttt{ConsultarPainel}) e as portas de saída
(\texttt{OcorrenciaRepository}, \texttt{EventoPublisher},
\texttt{GeocodificacaoGateway}, \texttt{ProtocoloGateway}); framework web,
persistência, mensageria e clientes externos são adaptadores que implementam
essas interfaces. Nenhuma classe do núcleo importa biblioteca de
infraestrutura \\
\addlinespace
Alternativas consideradas & \textbf{Arquitetura em Camadas} — descartada por não
impedir que a regra de negócio dependa do ORM e do framework, deixando sem
solução justamente o ponto crítico do sistema. \textbf{Clean Architecture} —
adequada e equivalente em efeito, descartada pelo número maior de camadas
prescritivas, cuja cerimônia não se paga neste porte. \textbf{Onion
Architecture} — próxima em princípio, descartada pela mesma razão, somada à
menor familiaridade da equipe com o seu vocabulário \\
\addlinespace
Consequências positivas & Troca de fornecedor externo passa a ser a escrita de
um adaptador, sem alteração de regra de negócio; testes de unidade do domínio
rodam sem banco, fila ou servidor, sustentando a base da pirâmide de testes;
novos canais de entrada não exigem reescrita do núcleo; a fronteira entre
negócio e tecnologia fica explícita no código e é verificável em revisão \\
\addlinespace
Consequências negativas & Mais arquivos e mais indireção do que uma organização
em camadas, com custo real para uma equipe de três integrantes; curva de
aprendizado da inversão de dependência, que exige disciplina em revisão de
código para não ser corrompida; risco de portas com implementador único, sem
ganho — mitigado pela regra de criar porta apenas onde houver volatilidade ou
necessidade de dublê em teste \\
\bottomrule
\end{tabularx}}
\fonte{elaborado pelos autores (2026), conforme o modelo de Nygard (2011).}
\end{table}

O Apêndice~A traz um segundo registro, não exigido pela atividade, sobre a
adoção de mensageria assíncrona — decisão que aparece no diagrama de containers
e que, por consistência, também convém documentar.

% =========================================================== 7
\section{ETAPA 6 — QUALIDADE E TESTES}
\label{sec:etapa6}

\subsection{Pirâmide de testes}

A Figura~\ref{fig:piramide} apresenta a estratégia de testes do NAD organizada
pela pirâmide. A regra de alocação é direta: o que é regra de negócio pura fica
na base; o que depende de recurso externo real fica no meio; e só sobe ao topo
o que precisa atravessar o sistema inteiro pela perspectiva do usuário.

\begin{figure}[H]
\centering
\caption{Pirâmide de testes do Neighborhood Action Dashboard}
\label{fig:piramide}
\includegraphics[width=\textwidth]{fig_piramide.png}
\fonte{elaborado pelos autores (2026), conforme Cohn (2009).}
\end{figure}

\subsection{Tipos de teste}

\begin{table}[H]
\centering
\caption{Níveis de teste e exemplos no Neighborhood Action Dashboard}
\label{tab:tipos}
{\small\singlespacing
\begin{tabularx}{\textwidth}{L{2.8cm} L{4.6cm} X}
\toprule
\cab{Nível} & \cab{Objetivo} & \cab{Exemplo no NAD} \\
\midrule
Testes Unitários & Verificar, de forma isolada, regras de negócio e funções
individuais & Política de priorização de uma demanda: categoria crítica com
reincidência no mesmo ponto deve resultar em prioridade alta. Validação dos
campos obrigatórios da ocorrência. Transição proibida de \emph{Aberta} para
\emph{Concluída} sem passar por \emph{Em execução} \\
\addlinespace
Testes de Integração & Verificar a comunicação entre componentes, serviços e
persistência & A API de ocorrências gravando e consultando no PostgreSQL, com a
consulta geoespacial que encontra demandas no raio de agrupamento; publicação e
consumo do evento \texttt{OcorrenciaRegistrada} na fila \\
\addlinespace
Testes End-to-End & Verificar fluxos completos sob a perspectiva do usuário &
O cidadão registra uma ocorrência e recebe o protocolo; o gestor prioriza,
atualiza a situação e encerra; o cidadão consulta o andamento e encontra o
histórico completo \\
\bottomrule
\end{tabularx}}
\fonte{elaborado pelos autores (2026).}
\end{table}

\subsection{Definition of Done}

A Tabela~\ref{tab:dod} define os critérios mínimos para considerar uma
funcionalidade pronta. A coluna \emph{Cumprido} é preenchida a cada
funcionalidade entregue; a coluna \emph{Evidência} registra o que comprova o
atendimento — sem evidência, o critério não é considerado cumprido.

\begin{table}[H]
\centering
\caption{Definition of Done do Neighborhood Action Dashboard}
\label{tab:dod}
{\small\singlespacing
\begin{tabularx}{\textwidth}{X C{1.8cm} L{5.4cm}}
\toprule
\cab{Critério mínimo} & \cab{Cumprido} & \cab{Evidência exigida} \\
\midrule
Requisito implementado e critérios de aceite validados & \vazio & Critérios de
aceite da história marcados e demonstrados \\
Testes unitários escritos e aprovados & \vazio & Execução verde no pipeline e
cobertura do pacote de domínio sem queda \\
Testes de integração aprovados & \vazio & Execução verde com banco e fila em
contêineres efêmeros \\
Fluxo principal validado em teste End-to-End & \vazio & Execução da jornada
correspondente registrada \\
Código revisado por outro integrante da equipe & \vazio & \emph{Pull request}
aprovado por um par \\
Decisão arquitetural registrada em ADR, quando aplicável & \vazio & ADR criado
ou atualizado no repositório \\
Sem defeitos críticos ou bloqueantes em aberto & \vazio & Quadro de defeitos sem
item de severidade 1 ou 2 associado \\
Documentação da funcionalidade atualizada & \vazio & Contrato OpenAPI e
anotações de uso revisados \\
\bottomrule
\end{tabularx}}
\fonte{elaborado pelos autores (2026), a partir do modelo proposto na GS2026.2.}
\end{table}

\subsection{Critério de cobertura}

A meta de cobertura é diferenciada por pacote: 80 \% de instruções e 70 \% de
ramos no núcleo de domínio, 60 \% nos adaptadores, e nenhuma meta para código
gerado e objetos de transporte. A regra que acompanha a meta importa mais do que
o número: a cobertura do domínio não pode cair de uma entrega para a seguinte.

Vale a ressalva, porém, de que \textbf{cobertura mede execução, não detecção}:
linhas percorridas por asserções fracas produzem a aparência de qualidade. Por
isso a Definition of Done exige critérios de aceite validados, e não apenas
testes verdes.

% =========================================================== 8
\section{CONSIDERAÇÕES FINAIS}

A proposta aqui apresentada parte de um diagnóstico simples: o problema do
Neighborhood Action Dashboard não é a falta de canais para o cidadão reclamar —
canais existem e se multiplicam. O problema é a ausência de um fluxo único,
padronizado e transparente, que transforme relato disperso em demanda
rastreável, priorizada por critério e verificável por quem reclamou.

A arquitetura escolhida responde ao traço mais marcante desse contexto: a
instabilidade das bordas. Canais de entrada se multiplicam e fornecedores
municipais mudam, enquanto a regra de negócio permanece. A Arquitetura
Hexagonal foi adotada exatamente para que a volatilidade fique confinada aos
adaptadores — e o documento registra também o que essa escolha custa, porque uma
decisão apresentada só pelos benefícios não é uma decisão documentada, é uma
defesa.

As visões em C4 tornam essa escolha verificável: o diagrama de contexto mostra
que quatro dos oito elementos externos são serviços substituíveis, e o de
containers mostra que todos eles são alcançados a partir de um único container,
pela via dos adaptadores. O cenário de registro de ocorrência, por sua vez,
mostra a arquitetura em funcionamento, inclusive naquilo que costuma ficar de
fora de uma proposta: o que acontece quando o dado é inválido, quando a
persistência falha e quando o problema já tinha sido reportado por outra pessoa.

Registre-se, por fim, o limite desta entrega. Trata-se de concepção
arquitetural, não de implementação: não há código, não há medição de desempenho
e os números de requisito não funcional são metas de projeto, não resultados.
A validação desses valores depende de protótipo executável e de uso real, que
são o passo seguinte natural deste trabalho.

% =========================================================== REFERÊNCIAS
\newpage
\phantomsection
\addcontentsline{toc}{section}{REFERÊNCIAS}
\section*{REFERÊNCIAS}
\vspace{0.6\baselineskip}

{\singlespacing\raggedright
\setlength{\parindent}{0pt}
\setlength{\parskip}{0.7\baselineskip}

ASSOCIAÇÃO BRASILEIRA DE NORMAS TÉCNICAS. \textbf{NBR 14724}: informação e
documentação: trabalhos acadêmicos: apresentação. Rio de Janeiro, 2011.

ASSOCIAÇÃO BRASILEIRA DE NORMAS TÉCNICAS. \textbf{NBR 6023}: informação e
documentação: referências: elaboração. Rio de Janeiro, 2018.

BRASIL. \textbf{Lei n.º 13.709, de 14 de agosto de 2018}. Lei Geral de Proteção
de Dados Pessoais (LGPD). Brasília: Presidência da República, 2018.

BROWN, Simon. \textbf{The C4 model for visualising software architecture}.
Disponível em: https://c4model.com. Acesso em: 6 nov.\ 2026.

COCKBURN, Alistair. \textbf{Hexagonal architecture}: ports and adapters. 2005.
Disponível em: https://alistair.cockburn.us/hexagonal-architecture. Acesso em:
6 nov.\ 2026.

COHN, Mike. \textbf{Succeeding with Agile}: software development using Scrum.
Boston: Addison-Wesley, 2009.

EVANS, Eric. \textbf{Domain-Driven Design}: tackling complexity in the heart of
software. Boston: Addison-Wesley, 2003.

FOWLER, Martin. \textbf{Refactoring}: improving the design of existing code.
2. ed. Boston: Addison-Wesley, 2019.

KRUCHTEN, Philippe. Architectural blueprints: the 4+1 view model of software
architecture. \textbf{IEEE Software}, v. 12, n. 6, p. 42-50, 1995.

MARTIN, Robert C. \textbf{Arquitetura limpa}: o guia do artesão para estrutura e
design de software. Rio de Janeiro: Alta Books, 2019.

NYGARD, Michael. \textbf{Documenting architecture decisions}. 2011. Cognitect
Blog. Disponível em:
https://www.cognitect.com/blog/2011/11/15/documenting-architecture-decisions.
Acesso em: 6 nov.\ 2026.

OBJECT MANAGEMENT GROUP. \textbf{OMG Unified Modeling Language (OMG UML)},
version 2.5.1. 2017. Especificação técnica formal/2017-12-05.

PALERMO, Jeffrey. \textbf{The Onion Architecture}. 2008. Disponível em:
https://jeffreypalermo.com/2008/07/the-onion-architecture-part-1. Acesso em:
6 nov.\ 2026.

SILVA, Leandro Ramos da. \textbf{GS2026.2 — Global Solution}: Neighborhood
Action Dashboard — CitizenTech para gestão colaborativa de demandas urbanas. São
Paulo: FIAP, 2026. Material de apoio da disciplina de Engenharia de Software.

WORLD WIDE WEB CONSORTIUM. \textbf{Web Content Accessibility Guidelines (WCAG)
2.1}. W3C Recommendation, 2018. Disponível em: https://www.w3.org/TR/WCAG21.
Acesso em: 6 nov.\ 2026.

}

% =========================================================== APÊNDICE A
\newpage
\phantomsection
\addcontentsline{toc}{section}{APÊNDICE A — ADR-002 (REGISTRO ADICIONAL)}
\section*{APÊNDICE A — ADR-002 (REGISTRO ADICIONAL)}
\vspace{0.6\baselineskip}

O registro a seguir não é exigido pela atividade, que pede um ADR. Ele é
incluído porque a decisão aparece no diagrama de containers e porque a prática
adotada pela equipe é documentar toda decisão com consequência estrutural.

\begin{table}[H]
\centering
\caption{ADR-002 — Processamento assíncrono de notificações e indicadores}
\label{tab:adr2}
{\small\singlespacing
\begin{tabularx}{\textwidth}{L{3.4cm} X}
\toprule
\cab{Status / Data} & \textbf{Aceita} \quad Data: 06/11/2026 \\
\midrule
Contexto & O registro de uma ocorrência precisa responder em menos de 2 s
(RNF01), mas dispara duas atividades que dependem de terceiros ou de agregação:
notificar o cidadão e atualizar os indicadores do painel. Executá-las dentro da
requisição significa que a lentidão ou a queda do serviço de notificação se
transforma em falha no registro — justamente a operação que não pode falhar \\
\addlinespace
Decisão & Publicar o evento \texttt{OcorrenciaRegistrada} em fila de mensagens e
tratar notificação e indicadores em \emph{workers} independentes, de modo que a
resposta ao cidadão dependa apenas da persistência \\
\addlinespace
Alternativas consideradas & Chamada síncrona dentro da requisição — descartada
por acoplar a disponibilidade do registro à de terceiros. Rotina periódica em
lote — descartada porque atrasaria a notificação em minutos e removeria o efeito
de transparência imediata que justifica a plataforma \\
\addlinespace
Consequências positivas & O registro continua funcionando com o serviço de
notificação indisponível; picos de demanda são absorvidos pela fila; novos
consumidores do evento podem ser acrescentados sem alterar o fluxo de registro \\
\addlinespace
Consequências negativas & Acrescenta um componente de infraestrutura a operar e
monitorar; introduz consistência eventual entre o registro e o painel de
indicadores; exige consumidores idempotentes, porque a entrega da mensagem pode
se repetir \\
\bottomrule
\end{tabularx}}
\fonte{elaborado pelos autores (2026), conforme o modelo de Nygard (2011).}
\end{table}

% =========================================================== APÊNDICE B
\newpage
\phantomsection
\addcontentsline{toc}{section}{APÊNDICE B — CHECKLIST FINAL}
\section*{APÊNDICE B — CHECKLIST FINAL}
\vspace{0.6\baselineskip}

\begin{table}[H]
\centering
\caption{Conferência dos itens obrigatórios do documento}
\label{tab:checklist}
{\small\singlespacing
\begin{tabularx}{\textwidth}{C{1.6cm} L{4.4cm} X}
\toprule
\cab{Item} & \cab{Exigência} & \cab{Onde está neste documento} \\
\midrule
\marcado & Capa & Página de rosto do documento \\
\marcado & Integrantes & Folha de rosto, com nome completo e RM \\
\marcado & Documento de Visão & Seção 2 — objetivo, problema, público-alvo,
benefícios, As-Is/To-Be, papéis e requisitos \\
\marcado & Arquitetura da Solução & Seção 3 — Arquitetura Hexagonal, com motivação,
benefícios, limitações e alternativas descartadas \\
\marcado & C4 ou 4+1 & Seção 4 — Modelo C4, diagramas de contexto e de containers \\
\marcado & Diagrama de Sequência & Seção 5 — cenário \emph{Registrar ocorrência},
com fluxo principal e três fluxos alternativos \\
\marcado & ADR & Seção 6 — ADR-001; Apêndice A traz o ADR-002, adicional \\
\marcado & Estratégia de Testes & Seção 7 — pirâmide, três níveis de teste e
Definition of Done \\
\marcado & Referências & Lista de referências conforme a ABNT NBR 6023 \\
\bottomrule
\end{tabularx}}
\fonte{elaborado pelos autores (2026), conforme o checklist final da GS2026.2.}
\end{table}

\end{document}
